\documentclass[11pt]{amsart}

\usepackage{amsmath,amssymb,amsthm,mathtools}
\usepackage[a4paper,margin=27mm]{geometry}
\usepackage[hidelinks]{hyperref}

\newtheorem{theorem}{Theorem}[section]
\newtheorem{proposition}[theorem]{Proposition}
\newtheorem{lemma}[theorem]{Lemma}
\theoremstyle{remark}
\newtheorem{remark}[theorem]{Remark}

\newcommand{\Z}{\mathbb Z}
\newcommand{\Q}{\mathbb Q}
\newcommand{\Sym}{\operatorname{Sym}}
\newcommand{\Hom}{\operatorname{Hom}}
\newcommand{\gm}{\gamma}
\newcommand{\aug}{\varpi}
\newcommand{\Rel}{\mathcal R}
\newcommand{\pr}{\operatorname{pr}}

\title[Odd-dimensional vanishing for finite metabelian Lie rings]
{A PBW and derived-functor proof of an odd-dimensional vanishing theorem\\
for finite metabelian Lie rings}
\author{}
\date{}

\begin{document}

\begin{abstract}
Let $L$ be a finite metabelian Lie ring of nilpotency class at most
$n+1$ and suppose that its last nonzero lower-central term
$\gamma_{n+1}(L)$ is a cyclic $2$-group.  We prove that
$\delta_{2n+1}(L)=0$.  The proof uses PBW collection and the ordinary
Koszul complexes for derived symmetric powers.  The intermediate PBW
packets need not be cycles; instead, a chain-level transgression kills
their contributions directly.  The last, quadratic, packet is a genuine
cycle and is annihilated by an integral Bockstein certificate.  At no
point is a homogeneous tail of a relation treated as a relation.
\end{abstract}

\maketitle

\section{Statement and preliminary reductions}

For a Lie ring $L$, let $U(L)$ be its universal enveloping algebra and
$\aug(L)$ its augmentation ideal.  Put
\[
 \delta_r(L)=L\cap\aug(L)^r,
 \qquad
 \gm_1(L)=L,
 \qquad
 \gm_{r+1}(L)=[\gm_r(L),L].
\]
All Lie rings in this paper are Lie rings over $\Z$.

\begin{theorem}\label{thm:main}
Let $n\geq1$, and let $L$ be a finite metabelian Lie ring such that
$\gm_{n+2}(L)=0$ and
\[
 C:=\gm_{n+1}(L)\cong\Z/2^e
\]
is cyclic.  Then
\[
                         \delta_{2n+1}(L)=0.
\]
\end{theorem}

Whenever $C=\langle z\rangle\cong\Z/q$ occurs below, we identify it
with $\frac1q\Z/\Z\subseteq\Q/\Z$ by $z\mapsto\frac1q+\Z$.

We use the following two standard facts.  For every Lie ring,
$\delta_3(L)=\gm_3(L)$.  For every metabelian Lie ring,
\begin{equation}\label{eq:PS}
                         2\delta_r(L)\subseteq\gm_r(L).
\end{equation}
The latter assertion is due to Passi and Sicking
\cite[Theorem~2.1]{PS}.

\begin{lemma}[cyclic last-layer reduction]\label{lem:cyclic-reduction}
Fix $s\geq2$.  Suppose Theorem~\ref{thm:main} is known with $s$ in
place of $n$ whenever $\gm_{s+1}$ is a cyclic $2$-group, and suppose
the theorem is already known in every class at most $s$.  Then it is
valid, with index $2s+1$, for every finite metabelian Lie ring of
class at most $s+1$.
\end{lemma}

\begin{proof}
Let $A$ be such a Lie ring and let $0\ne a\in\delta_{2s+1}(A)$.
Since $\gm_{2s+1}(A)=0$, \eqref{eq:PS} gives $2a=0$.  The image of
$a$ in $A/\gm_{s+1}(A)$ belongs to
$\delta_{2s+1}\subseteq\delta_{2s-1}=0$ by the hypothesis in lower
class.  Hence $a$ lies in the central group $\gm_{s+1}(A)$.  Its
$2$-primary part is a finite abelian $2$-group.
There is a homomorphism from that group to a cyclic $2$-group which is
nonzero on $a$: write the group as a direct sum of cyclic groups and
project onto a summand on which $a$ has nonzero coordinate.  Extend
the kernel by the odd-primary part.  This kernel is an ideal because
$\gm_{s+1}(A)$ is central.  In the resulting quotient the image of
$a$ is nonzero, still belongs to $\delta_{2s+1}$, and the last
lower-central term is cyclic of $2$-power order, contradicting the
hypothesis.
\end{proof}

\begin{proposition}[reduction to the top layer]\label{prop:top-reduction}
To prove Theorem~\ref{thm:main} it is enough to show that, for
$C=\langle z\rangle\cong\Z/q$, $q=2^e$,
\begin{equation}\label{eq:top-reduction}
 \zeta z\in\delta_{2n+1}(L)
 \quad\Longrightarrow\quad
 \zeta\equiv0\pmod q.
\end{equation}
\end{proposition}

\begin{proof}
We argue by induction on $n$.  If $n=1$, then
$\delta_3(L)=\gm_3(L)=0$.  Assume the assertion known for $n-1$.
For $n=2$, the equality $\delta_3(A)=\gm_3(A)$ gives
$\delta_3(A)=0$ for every $A$ of class at most two.  For $n>2$,
Lemma~\ref{lem:cyclic-reduction} gives
$\delta_{2n-1}(A)=0$ for every finite metabelian Lie ring $A$ of
class at most $n$.  If $a\in\delta_{2n+1}(L)$, then its image in
$L/C$ belongs to
\[
 \delta_{2n+1}(L/C)\subseteq\delta_{2n-1}(L/C)=0.
\]
Thus $a\in C$ and has the form $a=\zeta z$.  Now
\eqref{eq:top-reduction} gives $a=0$.
\end{proof}

\section{The graded free model and the PBW witness}

Choose a finite generating set of $L$ and write
\[
                         L=F/\Rel,
\]
where $F$ is the relatively free metabelian Lie ring of class
$n+1$ on a finitely generated free abelian group $X$.  The usual Hall
basis gives a grading by bracket weight
\[
 F=\bigoplus_{s=1}^{n+1}F_s,
 \qquad F_1=X,
 \qquad [F_i,F_j]\subseteq F_{i+j}.
\]
Every $F_s$ is free abelian.  Put $M=F_{\geq2}=[F,F]$.  Metabelianness
gives
\begin{equation}\label{eq:metabelian-grading}
 [F_i,F_j]=0\quad(i,j\geq2),
 \qquad
 [F_i,F_1]\subseteq F_{i+1}.
\end{equation}
In particular $M$ is a module over $\Sym(X)$ by
$m\cdot x=[m,x]$; Jacobi and $[M,M]=0$ make this action symmetric.

Let $U=U(F)$.  Order a homogeneous basis of $F$ first by weight and
then arbitrarily inside each weight.  PBW identifies $U$, as an
abelian group, with the free group on the corresponding ordered
monomials.  Give a PBW monomial total weight equal to the sum of the
weights of its Lie factors, and denote the weight-$w$ subgroup by
$U_w$.

\begin{lemma}\label{lem:augmentation-weight}
For every $r\geq1$,
\[
                         \aug(F)^r=\bigoplus_{w\geq r}U_w.
\]
\end{lemma}

\begin{proof}
Every element of $F_s$ is a bracket of length $s$ in $X$, hence lies
in $\aug(F)^s$.  Therefore every PBW monomial of total weight $w$
belongs to $\aug(F)^w$.  This proves the inclusion from right to left.
Conversely, $\aug(F)$ is generated by homogeneous elements of positive
weight, so a product of $r$ of its elements has weight at least $r$.
\end{proof}

Let $\mathfrak r$ be the two-sided ideal of $U$ generated by $\Rel$.
Since
\[
 x\rho=\rho x+[x,\rho]
 \qquad(x\in F,\ \rho\in\Rel)
\]
and $[x,\rho]\in\Rel$, one has
\begin{equation}\label{eq:left-ideal}
                         \mathfrak r=\Rel U.
\end{equation}
Fix $\widetilde z\in F_{n+1}$ mapping to $z$.  If
$\zeta z\in\delta_{2n+1}(L)$, then there is
\begin{equation}\label{eq:witness}
 \theta\in\Rel U,
 \qquad
 \theta-\zeta\widetilde z\in\aug(F)^{2n+1}.
\end{equation}
Write $\theta_w$ for the weight-$w$ part.
Lemma~\ref{lem:augmentation-weight} gives
\begin{equation}\label{eq:governing}
 \theta_w=0\quad(w\leq2n,\ w\ne n+1),
 \qquad
 \theta_{n+1}=\zeta\widetilde z.
\end{equation}
Moreover, within each fixed weight the PBW summands with different
numbers of Lie factors are independent.  Thus every multi-factor
PBW coefficient of $\theta$ of weight at most $2n$ is zero.
Since a one-factor PBW monomial has weight at most $n+1$, one also has
\begin{equation}\label{eq:primitive-witness}
                         \pr_1(\theta)=\zeta\widetilde z,
\end{equation}
where $\pr_1:U\to F$ denotes the one-factor PBW projection.

\section{Canonical presentations and Koszul complexes}

For $2\leq k\leq n$, put
\[
 U_k=\gm_k(L)/\gm_{k+1}(L),\qquad
 W_k=L/\gm_{k+1}(L),\qquad
 V_k=L/\gm_k(L).
\]
There is an exact sequence of additive groups
\begin{equation}\label{eq:tower-extension}
 E_k:\qquad 0\longrightarrow U_k\longrightarrow W_k
 \xrightarrow{\pi_k}V_k\longrightarrow0.
\end{equation}

There is a convenient presentation obtained directly from $F$.  Set
\[
 A_k=F_{\leq k},
 \qquad
 D_k=\pr_{\leq k}(\Rel)\subseteq A_k.
\]
Then
\begin{equation}\label{eq:canonical-resolution}
 0\longrightarrow D_k\xrightarrow{d_k}A_k
 \longrightarrow W_k\longrightarrow0
\end{equation}
is exact, where $d_k$ is inclusion.  Projection onto weights at most
$k-1$ gives a morphism from \eqref{eq:canonical-resolution} to the
analogous resolution
\[
 0\longrightarrow D_{k-1}\xrightarrow{d_{k-1}}A_{k-1}
 \longrightarrow V_k\longrightarrow0.
\]
Indeed, the kernel of $A_k\to L/\gm_{k+1}(L)$ is exactly
$\pr_{\leq k}(\Rel)=D_k$.

For a two-term free resolution $P=(P_1\xrightarrow dP_0)$, let
\begin{equation}\label{eq:koszul}
 K_m(P)_i=\Lambda^iP_1\otimes\Sym^{m-i}P_0
 \qquad(0\leq i\leq m)
\end{equation}
with the usual Koszul differential
\[
 \partial(a_1\wedge\cdots\wedge a_i\otimes u)
 =\sum_{j=1}^i(-1)^{i-j}
 a_1\wedge\cdots\widehat{a_j}\cdots\wedge a_i
 \otimes da_j\,u.
\]
Then $H_iK_m(P)=L_i\Sym^m(\operatorname{coker}d)$; see, for example,
\cite{BM}.

We use only these full complexes.  Write
\[
 K_m(k)=K_m(D_k\longrightarrow A_k).
\]
Projection onto weights at most $k-1$ induces a chain map
\[
 f_k:K_m(k)\longrightarrow K_m(k-1).
\]

\section{The higher characters and their chain contractions}

Put
\[
                         m_k=n+2-k.
\]
The bracket gives a well-defined symmetric map
\begin{equation}\label{eq:theta-k}
 \Theta_k:U_k\otimes\Sym^{m_k-1}(V_k)\longrightarrow C,
 \qquad
 \Theta_k(\bar u;x_1\cdots x_{m_k-1})
 =[u,x_1,\ldots,x_{m_k-1}].
\end{equation}
Indeed, changing $u$ by an element of $\gm_{k+1}$ produces a bracket
in $\gm_{n+2}=0$; changing one $x_i$ by an element of $\gm_k$ gives
a bracket of two elements of $L'$.  Symmetry follows from Jacobi and
$[L',L']=0$.

The extension cocycle of \eqref{eq:tower-extension} can be described
on the canonical resolutions.  If $a\in D_{k-1}$, choose
$\rho\in\Rel$ with $\pr_{\leq k-1}\rho=a$ and put
\[
 b_k(a)=-\rho_k\pmod{D_k\cap F_k}\in U_k.
\]
The result is independent of the choice of $\rho$.  On the
degree-one part of $K_{m_k}(D_{k-1}\to A_{k-1})$, define
\begin{equation}\label{eq:phi-k}
 \Phi_k(a;x_1\cdots x_{m_k-1})
 =-\Theta_k(b_k(a);\bar x_1\cdots\bar x_{m_k-1}).
\end{equation}
It is a cocycle: in the value on a Koszul boundary one argument is
$d_{k-1}a'$, whose image in $V_k$ is zero.  It is supported in total
bracket weight $n+1$.  Consequently it defines the natural character
\begin{equation}\label{eq:higher-eta}
 \eta_k:L_1\Sym^{m_k}(V_k)\longrightarrow\Q/\Z,
 \qquad \eta_k([c])=\Phi_k(c)
 \quad(2\leq k<n).
\end{equation}
This character depends only on the extension $E_k$ and $\Theta_k$.
Indeed, replacing the cocycle lift $b_k$ by
$b_k+\lambda d_{k-1}$ changes $\Phi_k$ by the coboundary of
\[
 S_k(x_1\cdots x_{m_k})
 =-\sum_i\Theta_k(\lambda x_i;
       \bar x_1\cdots\widehat{\bar x_i}\cdots\bar x_{m_k}).
\]
The usual comparison maps between free resolutions then show that
$\eta_k$ is natural.

For $k<n$ there is a chain-level null-homotopy after pullback.  Write
$w=w_Q+w_P$ according to $A_k=A_{k-1}\oplus F_k$, and denote the
images of these components in $V_k$ and $U_k$ by $\bar w_Q$ and
$\bar w_P$.  Define
\begin{equation}\label{eq:T-k}
 T_k(w_1\cdots w_{m_k})=
 \sum_{i=1}^{m_k}
 \Theta_k(\bar w_{i,P};
     \bar w_{1,Q}\cdots\widehat{\bar w_{i,Q}}\cdots
     \bar w_{m_k,Q}).
\end{equation}

\begin{lemma}[transgression]\label{lem:transgression}
For $2\leq k<n$,
\[
                         f_k^*\Phi_k=\delta T_k.
\]
In particular,
\begin{equation}\label{eq:chain-transgression}
 \Phi_k(f_kc)=T_k(\partial c)
 \qquad(c\in K_{m_k}(k)_1),
\end{equation}
and
\begin{equation}\label{eq:higher-vanishing}
 \eta_k\circ L_1\Sym^{m_k}(\pi_k)=0.
\end{equation}
\end{lemma}

\begin{proof}
Let $r\in D_k$ and $w_2,\ldots,w_{m_k}\in A_k$.  The
$F_k$-component of $d_kr$ is the negative of the cocycle tail
$b_k(\pr_{\leq k-1}r)$, modulo a relation in $U_k$.  In
$T_k((d_kr)w_2\cdots w_{m_k})$, the summand in which the first factor
occupies the $U_k$-slot is therefore
$f_k^*\Phi_k(r;w_2\cdots w_{m_k})$.  Every other summand
contains the image of $d_kr$ in $V_k$ and is zero.  This proves the
identity and hence \eqref{eq:chain-transgression}.  Evaluation on a
cycle gives \eqref{eq:higher-vanishing}.  Notice also that $T_k$ can
be nonzero only when its $F_k$-argument has weight $k$ and all other
arguments have weight one: if another argument has weight at least
two, \eqref{eq:metabelian-grading} makes the long bracket zero.  Thus
$T_k$ sees precisely total bracket weight
$k+(m_k-1)=n+1$.  This support statement is all that the PBW argument
will need.
\end{proof}

\section{The quadratic character}

We now treat $k=n$, so $m_n=2$.  Write
\[
 0\longrightarrow U\longrightarrow W\xrightarrow\pi V
 \longrightarrow0
\]
for $E_n$, and retain $C=\Z/q$.  Choose Smith free resolutions
$P_1\xrightarrow dP_0\to V$ and $Q_1\xrightarrow eQ_0\to U$, and a
lift $\widetilde b:P_1\to Q_0$ of the extension cocycle.  A resolution
of $W$ is
\begin{equation}\label{eq:quadratic-resolution}
 R_1=P_1\oplus Q_1\xrightarrow\partial R_0=P_0\oplus Q_0,
 \qquad
 \partial(a,s)=(da,es-\widetilde b(a)).
\end{equation}

Let $g:V\otimes U\to C$ be $g(x,u)=[u,x]$.  Define
\begin{equation}\label{eq:quadratic-phi}
 \varphi(a,x)=g(\bar x,\overline{\widetilde b(a)}).
\end{equation}
This is a cocycle on $K_2(P)$, since the image of $da'$ in $V$ is zero.
Choose an integral lift $\widehat\varphi$.  It may be chosen so that
\begin{equation}\label{eq:exterior-lift}
 \widehat\varphi(a,da')+\widehat\varphi(a',da)=0.
\end{equation}
Indeed, write $da_i=d_ix_i$, with $d_i\mid d_j$ for $i<j$, and lift
$x_i$ to $X_i\in L$.  The extension cocycle is represented by
$d_iX_i\in U$, so
\[
 \varphi(a_i,x_j)=d_i[X_i,X_j],
 \qquad
 \varphi(a_j,x_i)=-d_j[X_i,X_j]
 =-\frac{d_j}{d_i}\varphi(a_i,x_j)
\]
in $C$.  Choose $r_{ij}\in\Z$ with
$r_{ij}/q+\Z=\varphi(a_i,x_j)$, and set
\[
 \widehat\varphi(a_i,x_j)=r_{ij},
 \qquad
 \widehat\varphi(a_j,x_i)=-\frac{d_j}{d_i}r_{ij},
 \qquad
 \widehat\varphi(a_i,x_i)=0.
\]
This gives
\eqref{eq:exterior-lift} exactly over $\Z$.  The Koszul homology and the
resulting character are independent of the chosen free resolution.

Because $V$ is finite, $d\otimes\Q$ is an isomorphism and
$H_2K_2(P)=0$.  Universal coefficients give
\begin{equation}\label{eq:UCT}
 H^2\Hom(K_2(P),\Z)
 \cong\Hom(L_1\Sym^2(V),\Q/\Z).
\end{equation}
The integral cochain
\[
 h_{\widehat\varphi}(a\wedge a')
 =\frac{\widehat\varphi(a,da')}{q}
\]
is well defined because $da'$ represents zero in $V$, so the numerator
is divisible by $q$.  It
defines, by \eqref{eq:UCT}, a character
\begin{equation}\label{eq:eta}
                         \eta_n:L_1\Sym^2(V)\to\Q/\Z.
\end{equation}
If $\beta_S$ is the Bockstein associated with
$0\to\Z\xrightarrow q\Z\to C\to0$, then
\begin{equation}\label{eq:half}
                         \beta_S[\varphi]=2\eta_n.
\end{equation}
Indeed, on $a\wedge a'$ its integral representative is
\[
 \frac{\widehat\varphi(a,da')-\widehat\varphi(a',da)}q
 =2h_{\widehat\varphi}(a\wedge a')
\]
by \eqref{eq:exterior-lift}.
Thus $\eta_n$, rather than its double
$\beta_S[\varphi]=2\eta_n$, is the quadratic PBW character.

\begin{proposition}[canonical quadratic PBW block]
\label{prop:quadratic-block}
Suppose that \eqref{eq:quadratic-resolution} is realized by the graded
presentation $F\to L$ used above.  Choose Smith bases
\[
 da_i=d_ix_i,
 \qquad es_\alpha=e_\alpha y_\alpha,
 \qquad d_i\mid d_j\quad(i<j),
\]
put $B_i=\widetilde b(a_i)$, and order every $x_i$ before every
$y_\alpha$ in the PBW basis.  Choose full relations
\begin{equation}\label{eq:quadratic-full-relations}
 \rho_i=d_ix_i-B_i+t_i,
 \qquad
 \sigma_\alpha=e_\alpha y_\alpha+t'_\alpha,
 \qquad t_i,t'_\alpha\in F_{n+1}.
\end{equation}
Let
\[
 \chi=c+v+v'+w\in
 (P_1\otimes P_0)\oplus(P_1\otimes Q_0)
 \oplus(Q_1\otimes P_0)\oplus(Q_1\otimes Q_0)
\]
be a cycle.  It has a canonical placed realization obtained as follows:
write
\begin{align*}
 c&=\sum_{i<j}u_{ij}
 \left(\frac{d_j}{d_i}a_i\otimes x_j-a_j\otimes x_i\right),\\
 v&=\sum v_{i\alpha}a_i\otimes y_\alpha,
 &v'&=\sum v'_{\alpha i}s_\alpha\otimes x_i,
 &w&=\sum w_{\alpha\beta}s_\alpha\otimes y_\beta,
\end{align*}
and use the placed rows
\begin{equation}\label{eq:canonical-quadratic-rows}
 \begin{split}
 &u_{ij}\left(\frac{d_j}{d_i}\rho_i x_j-x_i\rho_j\right),
 \qquad v_{i\alpha}\rho_i y_\alpha,\\
 &v'_{\alpha i}x_i\sigma_\alpha,
 \qquad w_{\alpha\beta}\sigma_\alpha y_\beta.
 \end{split}
\end{equation}
If the remaining one-factor coefficient of this complete block in $C$
is $\kappa z$, then
\begin{equation}\label{eq:quadratic-PBW-block}
 \frac\kappa q+\Z
 =-\eta_n\bigl(L_1\Sym^2(\pi)[\chi]\bigr).
\end{equation}
Every placed two-factor block with vanishing two-factor coefficients is
carried to \eqref{eq:canonical-quadratic-rows} by integral relation
interchanges; their correction terms are full one-factor relations.
\end{proposition}

\begin{proof}
First, a cycle in the $P_1\otimes P_0$ summand has the displayed form.
Indeed, the coefficient of $x_i^2$ is $d_i$ times the diagonal
coefficient and hence forces that coefficient to be zero; the
$x_ix_j$ coefficient gives
$d_i c_{ij}+d_jc_{ji}=0$, and $d_i\mid d_j$.  The remaining components
of $\partial\chi=0$, sorted in
\[
 \Sym^2(P_0\oplus Q_0)=
 \Sym^2P_0\oplus(P_0\otimes Q_0)\oplus\Sym^2Q_0,
\]
are
\begin{align}
 (d\otimes1)c&=0,\label{eq:block-A}\\
 -(\widetilde b\otimes1)c
 +(d\otimes1)v+(e\otimes1)v'&=0,\label{eq:block-B}\\
 -\mu(\widetilde b\otimes1)v+(e\otimes1)w&=0.
 \label{eq:block-C}
\end{align}
Thus the two-factor PBW symbol of \eqref{eq:canonical-quadratic-rows}
is zero.  Conversely, an adjacent interchange
$x\rho=\rho x+[x,\rho]$ changes no symmetric two-factor symbol and
adds a full one-factor relation.  It therefore converts any placement
with this symbol, integrally, to the canonical one.

We now calculate its primitive term, rather than infer it from
\eqref{eq:block-A}--\eqref{eq:block-C}.  In the horizontal $i,j$ row,
the $P_0P_0$ products cancel before any interchange:
\[
 \frac{d_j}{d_i}d_ix_ix_j-x_id_jx_j=0.
\]
The $F_n$-tail of its first relation gives
\[
 -\frac{d_j}{d_i}B_ix_j
 =-\frac{d_j}{d_i}x_jB_i
  -\frac{d_j}{d_i}[B_i,x_j],
\]
whereas $x_iB_j$ in the second placement is already ordered.  Hence
the horizontal primitive coefficient is
\begin{equation}\label{eq:block-Z}
 \kappa\equiv-
 \sum_{i<j}u_{ij}\frac{d_j}{d_i}
 \widehat\varphi(a_i,x_j)\pmod q.
\end{equation}
There are no suppressed vertical corrections.  In $\rho_i y_\alpha$
the only possibly inverted tail product is $B_i y_\alpha$, and its
commutator is zero because both factors lie in the metabelian derived
ideal.  The products $x_i\sigma_\alpha$ are already ordered, and the
$\sigma_\alpha y_\beta$ products again have two derived factors.
Finally, the $t_i,t'_\alpha$ in
\eqref{eq:quadratic-full-relations} have weight $n+1$, so multiplying
them by a positive-weight factor cannot produce a primitive term in
weight $n+1$.  This proves \eqref{eq:block-Z} coefficient by
coefficient.

On the other hand,
\[
 \partial(a_i\wedge a_j)=d_i
 \left(\frac{d_j}{d_i}a_i\otimes x_j-a_j\otimes x_i\right),
 \qquad
 h_{\widehat\varphi}(a_i\wedge a_j)
 =\frac{d_j\widehat\varphi(a_i,x_j)}q.
\]
The universal-coefficient description \eqref{eq:UCT} gives
\[
 \eta_n([c])=
 \frac1q\sum_{i<j}u_{ij}\frac{d_j}{d_i}
 \widehat\varphi(a_i,x_j)+\Z.
\]
Projection of $[\chi]$ along $\pi$ is $[c]$; comparison with
\eqref{eq:block-Z} proves \eqref{eq:quadratic-PBW-block}.
\end{proof}

\subsection*{The certificate}

Put $A=\Hom(U,C)$ and $A^\vee=\Hom(A,C)\cong U/qU$, where the last
isomorphism is the evaluation map (it is an isomorphism on each cyclic
summand of $U$).  Let
\[
 \rho:P_0\to A,
 \quad \rho(x)=g(\bar x,-),
 \qquad
 b:P_1\to A^\vee,
 \quad b(a)=\overline{\widetilde b(a)}.
\]
Form the pushout
\[
 \mathcal E=(A^\vee\oplus P_0)/
 \langle(b(a),-da):a\in P_1\rangle.
\]
The map $\mathcal E\to V$, $[\chi,x]\mapsto\bar x$, is onto and its
kernel is generated by the finite group $A^\vee$; hence $\mathcal E$ is
finite.
Extend $\widehat\varphi$ rationally and define the alternating form
$\vartheta$ on $P_0\otimes\Q$ by
\[
 \vartheta(da,x)=\widehat\varphi(a,x)/q.
\]
Equation \eqref{eq:exterior-lift} makes it alternating.  On
$\mathcal E$ define
\[
 \Omega([\chi,x],[\chi',x'])
 =\chi(\rho x')-\chi'(\rho x)+\vartheta(x,x')\pmod\Z.
\]
It is well defined: substituting $(b(a),-da)$ in either variable
gives
$\varphi(a,x)-\widehat\varphi(a,x)/q=0$ in $\Q/\Z$.

Every alternating form on a finite abelian group with values in the
divisible group $\Q/\Z$ is the antisymmetrization of a bilinear form:
choose cyclic generators and place its matrix in one triangle.  Choose
such a $\Q/\Z$-valued form $H$ with
$H^{\mathsf T}-H=\Omega$.  Although $H$ need not be $C$-valued on all
of $\mathcal E\times\mathcal E$, its value is $C$-valued whenever one
argument lies in $i(A^\vee)$: every element of $A^\vee\cong U/qU$ is
killed by $q$, and bilinearity then kills the corresponding value by
$q$.
If
$i(\chi)=[\chi,0]$ and $j(x)=[0,x]$, define
\[
 \chi(\alpha x)=H(jx,i\chi),
 \qquad
 \lambda(\chi,\chi')=H(i\chi',i\chi).
\]
The preceding observation and the evaluation duality
$\Hom(A^\vee,C)\cong A$ show that $\alpha:P_0\to A$ is well defined.
Also $\lambda$ is $C$-valued and symmetric, and the relation
$jd=ib$ gives
\begin{equation}\label{eq:certificate-relation}
                         \alpha d=\lambda_*b.
\end{equation}
Indeed, for $\chi\in A^\vee$,
\[
 \chi(\alpha da)=H(jda,i\chi)=H(ib(a),i\chi)
 =\lambda(\chi,b(a)).
\]
Moreover, the cocycle
\[
                         \psi(a,x)=b(a)(\alpha x)
\]
satisfies
\begin{equation}\label{eq:certificate-bockstein}
                         \beta_S[\psi]=\eta_n.
\end{equation}
It is a cocycle because
\[
 \psi(a,da')-\psi(a',da)
 =\lambda(b(a),b(a'))-\lambda(b(a'),b(a))=0.
\]
To compute its Bockstein, lift
$(x,x')\mapsto H(jx,jx')$ to a rational bilinear form
$P$ on the free group $P_0$.  The restriction
$H(jP_0,jdP_1)=H(jP_0,ib(P_1))$ is $C$-valued, so every such lift
satisfies $qP(x,da)\in\Z$.  Adding an integral bilinear form in one
triangle, arrange $P^{\mathsf T}-P=\vartheta$ without changing that
integrality.  Then $\widehat\psi(a,x)=qP(x,da)$ is integral, reduces to
$\psi$, and
\[
 \widehat\psi(a,da')-\widehat\psi(a',da)
 =\widehat\varphi(a,da').
\]
After division by $q$, this is exactly
\eqref{eq:certificate-bockstein}.

\begin{proposition}[quadratic vanishing]\label{prop:quadratic-vanishing}
\[
                         \eta_n\circ L_1\Sym^2(\pi)=0.
\]
\end{proposition}

\begin{proof}
On $R_0=P_0\oplus Q_0$ define the symmetric $C$-valued form
\[
 \begin{split}
 T((x,y),(x',y'))={}&-\alpha(x)(y'_U)
                     -\alpha(x')(y_U)\\
                    &-\lambda(\bar y,\bar y'),
 \end{split}
\]
where $y_U$ denotes the image of $y$ in $U$ and $\bar y$ its image in
$U/qU$.  Using
\eqref{eq:quadratic-resolution} and
\eqref{eq:certificate-relation}, one obtains
\[
 \begin{aligned}
 T(\partial(a,s),(x,y))
 &=-\lambda(b(a),\bar y)+b(a)(\alpha x)
   +\lambda(b(a),\bar y)\\
 &=\psi(a,x).
 \end{aligned}
\]
Thus the pullback of $[\psi]$ to $K_2(R)$ is zero.  Naturality of the
Bockstein and \eqref{eq:certificate-bockstein} give
\[
 \pi^*\eta_n=\beta_S(\pi^*[\psi])=0,
\]
which is the assertion under \eqref{eq:UCT}.
\end{proof}


\section{The exact two-filtered row calculation}

The assembly must keep a relation whole while its homogeneous pieces are
being followed.  The bookkeeping device which does this is the relative
Chevalley--Eilenberg row complex
\begin{equation}\label{eq:CE-row-complex}
 \mathcal C_i=\Lambda^i\Rel\otimes U(F),\qquad
 \mathcal C_0=U(F),
\end{equation}
with differential
\begin{align}
 d(\rho_1\wedge\cdots\wedge\rho_i\otimes u)
={}&\sum_{a=1}^i(-1)^{i-a}
 \rho_1\wedge\cdots\widehat\rho_a\cdots\wedge\rho_i
 \otimes\rho_a u \notag\\
 &+\sum_{a<b}(-1)^{a+b}
 [\rho_a,\rho_b]\wedge
 \rho_1\wedge\cdots\widehat\rho_a\cdots
 \widehat\rho_b\cdots\wedge\rho_i\otimes u.
                                                        \label{eq:CE-differential}
\end{align}
It satisfies $d^2=0$, all relation entries in it are full relations, and
\begin{equation}\label{eq:CE-image}
 d(\rho\otimes u)=\rho u,\qquad
 \operatorname{im}(d:\mathcal C_1\to\mathcal C_0)=\Rel U.
\end{equation}

Give a tensor in $\mathcal C_i$ total factor number $i+t$ when its
$U(F)$-entry is a $t$-factor PBW monomial.  The factor-preserving symbol
of \eqref{eq:CE-differential}, in total factor number $p$, is the full
Koszul complex
\begin{equation}\label{eq:CE-Koszul-symbol}
 K_p(\Rel\hookrightarrow F)_i
 =\Lambda^i\Rel\otimes\Sym^{p-i}F.                    
\end{equation}
The lower-factor components are generated, integrally, by
\begin{equation}\label{eq:CE-transfers}
 \begin{split}
 [\cdots|v|u|\cdots]
  &=[\cdots|u|v|\cdots]+[\cdots|[v,u]|\cdots],\\
 [\cdots|v|\rho|\cdots]
  &=[\cdots|\rho|v|\cdots]+[\cdots|[v,\rho]|\cdots].
 \end{split}
\end{equation}
The correction in the second line again has a full relation entry.

To make the collection algorithm integral, filter $\Rel$ by
$\Rel\cap F_{\geq s}$.  The leading-component map embeds
\[
 (\Rel\cap F_{\geq s})/(\Rel\cap F_{\geq s+1})
 \lhook\joinrel\longrightarrow F_s.
\]
Its source is free abelian.  Choose bases successively and lift them;
this gives a triangular basis of $\Rel$ in which a basis relation of
leading weight $s$ has the form
\begin{equation}\label{eq:triangular-relation}
                         \rho=\rho_s+\rho_{>s}.
\end{equation}
A placed $p$-row is a word with $p-1$ homogeneous ordinary entries and
one distinguished basis relation.  During collection the relation in
the row always remains the full $\rho$.  Its active weight edge is the
identity
\begin{equation}\label{eq:active-weight-edge}
 \pr_{\geq s}\rho=\rho_s+\pr_{\geq s+1}\rho;
\end{equation}
the two terms on the right are recorded as components of the quotient
tower, not inserted into the row module as relations.

Iterating the first transfer in \eqref{eq:CE-transfers} gives, for
$B\in F_{\geq2}$ and $x_i\in F_1$, the literal integral identity
\begin{equation}\label{eq:subset-collection}
 Bx_1\cdots x_r
 =\sum_{I\subseteq\{1,\ldots,r\}}
   x_{I^c}[B;x_I],
\end{equation}
where both products and the teeth of the left-normed bracket retain
their original order.  This follows by induction on $r$ from
$Bx=xB+[B,x]$.  In particular the unique one-factor summand is the
$I=\{1,\ldots,r\}$ term, with coefficient $+1$.  Applying
\eqref{eq:subset-collection} simultaneously to all homogeneous
components in \eqref{eq:active-weight-edge} shows why the proper-subset
terms recombine into the full commutator rows required by
\eqref{eq:CE-transfers}.

For the second direction use the tower of quotient maps, rather than
splitting a tail off as a new relation.  The maps
\[
 \Rel\longrightarrow D_k,\quad \rho\longmapsto\pr_{\leq k}\rho,
 \qquad F\longrightarrow A_k
\]
take the factor-$p$ symbol to $K_p(k)$.  The passage from $k$ to $k-1$
is $f_k$, and
\begin{equation}\label{eq:successive-weight-differential}
 \pr_{\leq k}\rho=\pr_{\leq k-1}\rho+\rho_k
\end{equation}
is an equality between two components of the same full-row
differential.  In particular, $\rho_k$ is never asserted to belong to
$\Rel$.

Formally, put $A_{n+1}=F$, $D_{n+1}=\Rel$, and $A_0=D_0=0$.
For fixed $p$, the kernels of the surjections
\[
 K_p(n+1)\longrightarrow K_p(k)
\]
form a finite filtration of the actual Koszul symbol complex
$K_p(\Rel\hookrightarrow F)$; its successive quotient maps are the
$f_k$.  Thus the weight direction below is a filtration of complexes,
not a decomposition of $\rho$ inside $\Rel$.

We now record the finite two-filtered calculation which replaces the
tail-ledger assertion.

\begin{lemma}[closed-square row lemma]\label{lem:closed-square}
Let $\Theta\in\Rel U$ have zero multi-factor PBW coefficient in every
weight at most $2n$, and suppose that its one-factor part is homogeneous
of weight $n+1$.  Exact collection by \eqref{eq:CE-transfers} produces
chains
\begin{equation}\label{eq:packet-chains}
 \chi_k\in K_{m_k}(k)_1\quad(2\leq k<n),
 \qquad \chi_n\in Z_1K_2(n),
\end{equation}
with the following properties.

First, $\partial\chi_k$ is the complete factor-$m_k$ symbol of
$\Theta$ after projection to $\Sym^{m_k}A_k$.  Let
$Q_n(\chi_n)\in C$ denote the oriented primitive coefficient of the
placed factor-two block.  Define the terminal discrepancy, without any
claim that it is a relation tail, by
\begin{equation}\label{eq:terminal-definition}
 \varepsilon:=\iota_C(\pr_{1,n+1}\Theta)
 -\sum_{k=2}^{n-1}\Phi_k(f_k\chi_k)-Q_n(\chi_n),
\end{equation}
where $\iota_C:F_{n+1}\to C$ is the quotient map.  Then
\begin{equation}\label{eq:terminal-discrepancy}
 \varepsilon=-\sum_{k=2}^{n-1}T_k(\partial\chi_k).
\end{equation}
Consequently $\varepsilon=0$.
\end{lemma}

\begin{proof}
Start with any $x\in\mathcal C_1$ satisfying $dx=\Theta$ and expand its
relation entries in the triangular basis
\eqref{eq:triangular-relation}.  Process
factor number in descending order.  At factor number $p\geq3$, order
the ordinary entries and put the distinguished relation at the far
left; at $p=2$ retain the oriented Smith placements used in
Proposition~\ref{prop:quadratic-block}.  Expand a relation into
homogeneous components only to decide the PBW interchanges, and apply
the second identity in \eqref{eq:CE-transfers} simultaneously to the
full relation.  The component not yet reached is therefore a vertical
edge of the quotient tower, not a new relation row.  Append every
full-relation correction to the $(p-1)$-row list.  Within a fixed
factor level, an ordinary transfer lowers the PBW inversion number, a
relation transfer lowers the distance from its prescribed position,
	and a quotient-tower truncation lowers its mark.  Hence the
	lexicographic measure
	\[
	 (p,\ \text{unresolved adjacent pairs},\ k)
	\]
	strictly decreases at the step being processed.  On the weight
	truncation in the hypothesis it is finite, so the algorithm terminates.
Modulo the next factor level the differential is
\eqref{eq:CE-Koszul-symbol}.  Therefore, at $p=m_k$, replacing the full
relation by its image in $D_k$, projecting the other entries to $A_k$,
and forgetting placement gives a chain $\chi_k$ whose Koszul boundary is
exactly the complete projected $p$-factor symbol of $\Theta$.  This
proves the first assertion with integral coefficients.  When $k=n$,
every monomial in $\Sym^2A_n$ has total bracket weight at most $2n$;
the hypothesis on $\Theta$ therefore gives $\partial\chi_n=0$.

	It remains to compute the last exposed edge of this finite collection.
	We first fix occurrence presentations of all chains in
	\eqref{eq:packet-chains}.  Thus, after expanding ordinary entries into
	the chosen homogeneous PBW basis but before combining equal summands,
	write
	\begin{equation}\label{eq:indexed-packet}
	 \chi_k=\sum_{\alpha\in J_k}c_\alpha
	 (a_\alpha;x_{\alpha,1}\cdots x_{\alpha,m_k-1}),
	 \qquad a_\alpha=\pr_{\leq k}\rho_\alpha.
	\end{equation}
	Here $\rho_\alpha$ is the full relation carried by the placed row from
	which the summand arose, $c_\alpha\in\Z$ includes its sign, and two
	equal algebraic summands with different ancestry remain different
	elements of $J_k$.  Formula \eqref{eq:indexed-packet}, rather than an
	arbitrary expression of $\chi_k$ after like terms have been combined,
	is what is meant below by its occurrence presentation.

	For $2\leq k<n$, let $\Delta_k\subseteq J_k$ consist of the indices for
	which every $x_{\alpha,j}$ has weight one and the $F_k$-component
	$\rho_{\alpha,k}$ is the distinguished $M$-input.  Put
	$a^-_\alpha=\pr_{\leq k-1}\rho_\alpha$.  Since
	$b_k(a^-_\alpha)=-\rho_{\alpha,k}$ in $U_k$, linearity and the support
	of $\Theta_k$ give the equality of indexed sums
	\begin{equation}\label{eq:indexed-horizontal-read}
	 \Phi_k(f_k\chi_k)=
	 \sum_{\alpha\in\Delta_k}c_\alpha
	 [\rho_{\alpha,k},x_{\alpha,1},\ldots,
	                         x_{\alpha,m_k-1}].
	\end{equation}
	The brackets on the right of \eqref{eq:indexed-horizontal-read} are
	understood after passage from $F_{n+1}$ to $C$.
	This is also the exact primitive descendant of the homogeneous child of
	the truncation at $k$: by \eqref{eq:subset-collection} that child has
	one, and only one, one-factor summand, and its coefficient is $+1$.

	There is an equally explicit upper read.  Applying $T_k$ to the
	boundary of the occurrence indexed by $\alpha$ gives
	\[
	 T_k\bigl(a_\alpha x_{\alpha,1}\cdots
	                   x_{\alpha,m_k-1}\bigr).
	\]
	In \eqref{eq:T-k}, a term in which an ordinary factor supplies the
	$F_k$-slot also contains the image of
	$a^-_\alpha=d_{k-1}(f_k a_\alpha)$ in $V_k$, and is therefore zero.
	The term in which $a_\alpha$ supplies that slot is supported only when all
	the ordinary factors have weight one.  Consequently the formal supported
	summands read by $T_k(\partial\chi_k)$ (retained even
	when their value in $C$ is zero) are in coefficient-
	preserving bijection with $\Delta_k$, the bijection being
	\begin{equation}\label{eq:indexed-vertical-bijection}
	 \alpha\longmapsto
	 T_k\bigl(a_\alpha x_{\alpha,1}\cdots
	                   x_{\alpha,m_k-1}\bigr),
	\end{equation}
	and
	\begin{equation}\label{eq:indexed-vertical-read}
	 T_k(\partial\chi_k)=
	 \sum_{\alpha\in\Delta_k}c_\alpha
	\Theta_k(\bar\rho_{\alpha,k};
	       \bar x_{\alpha,1}\cdots\bar x_{\alpha,m_k-1}).
	\end{equation}
	Conversely, a formal supported term in the left side of
	\eqref{eq:indexed-vertical-read} has a unique index $\alpha$ before like
	terms are combined.  Its $F_k$-slot cannot be an ordinary slot by the
	preceding zero argument, so it is the relation slot and all remaining
	entries have weight one; hence $\alpha\in\Delta_k$.  This proves
	surjectivity as well as injectivity of
	\eqref{eq:indexed-vertical-bijection}.
	Repeated equal values cause no loss here: the index $\alpha$ and its
	coefficient are retained on both sides.  Equations
	\eqref{eq:indexed-horizontal-read} and
	\eqref{eq:indexed-vertical-read} are the occurrence-level form of
	\eqref{eq:chain-transgression}; in particular they prove both directions
	of the asserted correspondence, not only the construction of an upper
	term from a lower one.

	We now define the global ledger used to connect these local reads.  A
	$k$-marked row is a placed word in which the marked entry
	$\langle\rho,k\rangle$ evaluates as $\pr_{\leq k}\rho$; put
	$\langle\rho,n+1\rangle=\rho$ and
	$\langle\rho,0\rangle=0$.  Besides the first transfer in
	\eqref{eq:CE-transfers}, the rules involving a mark are
	\begin{align}
	 [\cdots|v|\langle\rho,k\rangle|\cdots]
	 &= [\cdots|\langle\rho,k\rangle|v|\cdots]
	   +[\cdots|\langle[v,\rho],k+h\rangle|\cdots],
	                                                        \label{eq:marked-transfer}\\
	 [\cdots|\langle\rho,k\rangle|\cdots]
	 &= [\cdots|\langle\rho,k-1\rangle|\cdots]
	   +[\cdots|\rho_k|\cdots]
	                                                        \label{eq:marked-truncation}
	\end{align}
	for homogeneous $v\in F_h$, with an index above $n+1$ read as
	$n+1$.  The second rule is literal.  The first is literal because
	\begin{equation}\label{eq:transfer-truncation-square}
	 \pr_{\leq k+h}[v,\rho]=[v,\pr_{\leq k}\rho].
	\end{equation}
	In particular, its correction is marked by the full relation
	$[v,\rho]\in\Rel$.

	Give every initial row a distinct finite-word label.  An output of a
	replacement receives its parent's word followed by the output number;
	a basis expansion appends the basis index as well.  Coefficients are
	multiplied along ancestry.  If $I$ is the initial signed occurrence list
	and $F_j$ is the list of outputs of the first $j$ replacements not yet
	used as inputs, then, in the free abelian group on these raw labels,
	\begin{equation}\label{eq:prefix-ledger}
	 I-F_j=\sum_{i=1}^{j}\partial_{\rm led}s_i,
	 \qquad
	 \partial_{\rm led}s_i=(\hbox{input})-(\hbox{signed outputs}).
	\end{equation}
	This follows by induction because the next step removes exactly its
	input and inserts exactly its outputs.  At termination it gives
	\begin{equation}\label{eq:complete-ledger}
	 I-F_N=\sum_{i=1}^{N}\partial_{\rm led}s_i.             
	\end{equation}
	This one-dimensional identity will be used only to show that no output
	occurrence is discarded.  The cancellation between two different
	orders of processing requires a separate incidence construction, which
	we give next.

	Fix once and for all the following deterministic choices: use the
	leftmost unresolved adjacent pair, expand in increasing basis order,
	and lower a mark by one.  On a signed factor-$p$ list let $V$ mean the
	complete deterministic pass which puts its factor-$p$ principal outputs
	in the prescribed order and returns every transfer correction, with its
	label, to the factor-$(p-1)$ list.  Let $H$ mean one marked truncation
	step, applied linearly to the whole list.  Thus $V$ and $H$ are
	homomorphisms of free abelian groups on occurrence lists; each fixes an
	output on which its indicated operation is unavailable.  A
	\emph{comparison cell} based
	at a list $R$ on which both operations are available is the formal
	square having the following four oriented incidence lists:
	\begin{equation}\label{eq:four-incidences}
	 \begin{array}{c|c}
	 \text{top}&h(R):=R-HR\\
	 \text{bottom}&-h(VR):=-VR+HVR\\
	 \text{right}&-v(R):=-R+VR\\
	 \text{left}&v(HR):=HR-VHR,
	 \end{array}
	\end{equation}
	where $v(S)=S-VS$.  Thus its oriented boundary is
	$HVR-VHR$.  This definition allows $H$ or $V$ to have several outputs;
	an ``incidence'' is the displayed coefficient-weighted list, not a
	single unlabelled value.

	Raw ancestry words along the two routes are normally different.  We do
	not identify them.  Instead, expand both route lists in the fixed placed
	basis and pair coefficient copies having the same initial prefix and the
	same resulting placed basis word.  The equality needed for this pairing
	is equality after evaluation and canonical basis expansion.  It follows
	from disjoint commutation, from ordinary Jacobi, and, when a relation is
	involved, from
	\begin{equation}\label{eq:CE-relation-Jacobi}
	 [u,[v,\rho]]-[v,[u,\rho]]=[[u,v],\rho].
	\end{equation}
	The remaining mixed case is exactly
	\eqref{eq:transfer-truncation-square}.  These cases exhaust the local
	peaks in the deterministic pass: an elementary step is either an
	adjacent transfer or a one-step truncation.  Induction on the number of
	unresolved adjacent pairs therefore shows, with the inherited integer
	coefficients, that $HVR=VHR$.  Hence every comparison cell has zero
	evaluated boundary.
	Notice the precise level of this assertion: different route labels are
	not equal as raw occurrences; their canonically expanded incidence
	\emph{lists} are coefficientwise equal.

	First form the \emph{uncut global trace} by starting with every initial
	label, adjoining its comparison cell whenever both $H$ and $V$ are
	available, and repeating this for every output list.  An incidence at
	the $V$-successor of one cell is, by definition, the top incidence of
	the adjacent cell; an incidence at the $H$-successor is the corresponding
	right incidence.  These are the only gluings.  Now cut this finite trace
	along the incidences in Rules (2) and (3) below and retain the side
	containing the factor-one frontier.  Equivalently, a cell is retained if
	it can be joined to a factor-one terminal occurrence without crossing a
	designated cut.  At a cut the incidence on that side is left unglued;
	the occurrence on the other side remains in the one-dimensional ledger
	\eqref{eq:complete-ledger}, so it has not been discarded.  Thus every internal
	incidence is contained in exactly two adjacent cells and occurs with
	opposite signs in \eqref{eq:four-incidences}.  Branching causes no
	ambiguity: it is expanded before gluing, and each coefficient copy has
	the unique initial prefix, operation, output number, and basis index
	specified above.  The critical-pair identities prove that the two lists
	being glued have equal evaluations; they do not manufacture a second
	ancestry for either raw label.

	The construction is finite because a correction lowers factor number,
	the principal output of a transfer lowers the number of unresolved
	adjacent pairs, and a truncation lowers the mark.  It is saturated by
	definition: after adjoining a cell, every nonsilent successor on which
	both operations remain available is adjoined in turn.  Equivalently,
	induction on the lexicographic triple
	\[
	 (p,\ \text{unresolved adjacent pairs},\ k)
	\]
	shows that every coefficient copy is either glued to its unique next
	incidence or reaches one of the stopping rules below.

	Here are the stopping rules; they are part of the definition of the
	trace, rather than an inference made after the calculation.
	\begin{enumerate}
	\item At factor number at least three, every unresolved ordinary or
	marked adjacency is processed.  Its principal output remains at that
	factor number and its correction is enqueued at the next lower factor
	number.  Once the principal output is canonical, its factor-$p$ product
	is retained as a vertical incidence; Rule (3) records it when it is on
	the support diagonal, and Rule (5) treats every other target read.
	\item At factor number two, all placements are retained until both
	entries have been projected to $A_n$.  The complete projected list is
	then put into the canonical Smith placements of
	Proposition~\ref{prop:quadratic-block}; every one-factor correction from
	those interchanges is enqueued in the one-dimensional ledger.  The
	common factor-two incidence is left unglued on the retained side and is
	the factor-two cut.  A component killed by the
	projection is recorded as silent only after the weight check below.
	\item Whenever $p=m_k$ and the mark crosses $k$, for
	$2\leq k<n$, record the homogeneous child and the corresponding
	factor-$p$ product incidence as the two indexed reads
	\eqref{eq:indexed-horizontal-read} and
	\eqref{eq:indexed-vertical-read}.  Leave those two distinct incidences
	unglued on the retained side, and follow each successor only in the
	one-dimensional ledger.  Thus every child remains recorded, while each
	diagonal boundary incidence occurs exactly once in the two-dimensional
	trace.
	\item At factor number one, all ordinary bracket trees are put into Hall
	form, and every remaining mark is truncated successively down to zero.
	Each homogeneous output is retained until its one-factor,
	weight-$(n+1)$ image in $C$ has been evaluated.
	\item An occurrence is declared silent only if it is zero in the placed
	basis, has a spectator after the one-factor projection, has total weight
	different from $n+1$, has two inputs in $M$, or is one of the full
	commutator relations identified below.
	\end{enumerate}
	These rules are exhaustive syntactically: a nonterminal row has factor
	number at least three and an available transfer, factor number two and
	the prescribed Smith processing, or factor number one and an available
	Hall or truncation step.  A canonical factor-$p$ product with $p\geq3$
	is the vertical incidence specified in Rule (1).  Thus there is no other
	kind of unprocessed terminal row.

	We emphasize what is cut in Rule (3).  Expand each of the four incidence
	lists in \eqref{eq:four-incidences} into its coefficient copies.  In the
	horizontal list $h(R)$, the parent and the lower-mark child are glued on
	as before; only the copy of the homogeneous child $\rho_k$ is exposed,
	with the minus sign in
	$R-(R^-+R_k)$.  The horizontal read of $R_k$ means its complete
	factor-lowering expansion followed by the one-factor target projection.
	Equation \eqref{eq:subset-collection} says that this read has the single
	term displayed in \eqref{eq:indexed-horizontal-read}; all proper-subset
	terms retain a spectator and have zero one-factor read.  On the vertical
	list, only the factor-$m_k$ product copy selected by the support of $T_k$
	is exposed, with the opposite boundary orientation, and its read is by
	definition \eqref{eq:indexed-vertical-bijection}.  Every other copy in
	those two incidence lists is still
	glued to its adjacent cell or is certified silent by Rule (5).  Thus a
	``diagonal incidence'' below means precisely one of these selected
	coefficient copies, not the whole replacement ledger $h(R)$ or $v(R)$.
	The same convention at the factor-two cut exposes exactly the copies in
	$B^{(n)}$; all Smith correction copies continue only in the
	one-dimensional ledger.

	We next verify the factor-two boundary without suppressing
	multiplicities.  Let $B$ be the signed list of factor-two occurrences
	just before the Smith processing.  An occurrence
	$c[\langle\rho,n+1\rangle|w]$ (or its opposite placement) contributes
	\begin{equation}\label{eq:factor-two-occurrence-map}
	 c\,(\pr_{\leq n}\rho;\pr_{\leq n}w)
	 \in D_n\otimes A_n
	\end{equation}
	to the occurrence presentation of $\chi_n$.  More precisely, let
	$B^{(n)}$ be the list obtained from $B$ by applying
	\eqref{eq:factor-two-occurrence-map}, deleting zero outputs, and
	expanding both entries in the fixed bases.  The definition of the
	occurrence presentation gives a bijection, including signs and
	multiplicities, between the coefficient copies in $B^{(n)}$ and its
	indexed summands.  Forgetting placement gives
	$\chi_n$, and taking the symmetric product gives $\partial\chi_n$.
	The latter is zero because it is precisely the projected factor-two PBW
	symbol of $\Theta$.

	No target-relevant correction is lost in
	\eqref{eq:factor-two-occurrence-map}.  Indeed, a later one-factor
	descendant is the bracket of the two homogeneous inputs and preserves
	their total weight.  If that weight is $n+1$, positivity gives
	$1\leq\operatorname{wt}(u),\operatorname{wt}(v)\leq n$, so both inputs
	survive projection to $A_n$.  Conversely, every summand of the
	occurrence presentation of $\chi_n$ comes from its unique element of
	$B^{(n)}$ by construction.  The Smith interchanges are literal integral
	identities and every correction keeps its coefficient and receives its
	output label.  Proposition~\ref{prop:quadratic-block} therefore reads
	the primitive contribution of exactly this list, with its multiplicity,
	as $Q_n(\chi_n)$; the factor-direction boundary orientation supplies
	the minus sign used below.

	It remains to determine which other stopping incidences can survive the
	target projection.  Every one-factor descendant is a bracket tree in
	the homogeneous inputs of its ancestor.  Expanding it in the fixed Hall
	basis is integral.  A nonzero metabelian Hall term has at most one input
	in $M=F_{\geq2}$, and the following alternatives partition the terminal
	coefficient copies.
	\begin{itemize}
	\item If the unique $M$-input is the relation component $\rho_k$, all
	the other inputs are in $F_1$.  A term descended from a factor-$p$ row
	has weight $k+p-1$; hence target weight forces $p=m_k$.  The cases
	$2\leq k<n$ are precisely the indices in $\Delta_k$, $k=n$ is precisely
	the projected factor-two list $B^{(n)}$, $k=n+1$ lies on the outside
	factor-one edge,
	and $k=1$ is covered by the full-commutator case below.
	\item If the unique $M$-input is ordinary, every relation component of
	weight at least two brackets to zero with it.  Replacing the weight-one
	component by the full relation therefore leaves the Hall term unchanged;
	it is a full relation commutator and maps to zero in $C$.  A term with
	two $M$-inputs is zero by \eqref{eq:metabelian-grading}.
	\item If every input has weight one, target weight forces $p=n+1$.
	Replacing the weight-one relation component by the full relation adds
	only terms of weight at least $n+2$, which vanish in $F$; the term is
	again a full relation commutator.  A spectator is killed by the
	one-factor projection, and $p>n+1$ has weight greater than the target.
	\end{itemize}
	This is exhaustive for occurrences, not just values: Rules (1)--(5)
	leave as nonsilent candidates only a one-factor Hall term, a retained factor-two term, or a
	diagonal incidence; the displayed alternatives partition the first of
	these by the number and source of its $M$-inputs.  The ancestry and basis
	indices then select exactly one alternative, even when two alternatives
	happen to evaluate to the same element of $F$.

	There is one read map throughout this boundary calculation.  For a
	placed incidence list $S$, let
	\[
	 \mathcal P(S)=\iota_C\bigl(\pr_{1,n+1}
	        (\text{the complete deterministic PBW collection of }S)\bigr).
	\]
	It is a homomorphism on signed lists and is unchanged by every gluing,
	because glued lists have the same canonical PBW expansion.  On a
	diagonal horizontal copy, \eqref{eq:subset-collection} computes
	$\mathcal P$ as \eqref{eq:indexed-horizontal-read}; on the corresponding
	vertical product copy, \eqref{eq:T-k} and
	\eqref{eq:indexed-vertical-read} compute the same map as $T_k$.  On the
	factor-two cut Proposition~\ref{prop:quadratic-block} computes it as
	$Q_n$, and on the occurrences in Rule (5) it is zero.  Thus the four
	entries below are evaluations of one map $\mathcal P$ on four kinds of
	incidences, not unrelated functionals attached after the gluing.

	Let $\mathcal G$ be the resulting finite multiset of comparison cells,
	with coefficient copies after canonical expansion.  The definition of
	the gluing gives the literal finite Stokes identity
	\begin{equation}\label{eq:finite-stokes}
	 \partial_{\rm ext}\mathcal G
	 =\sum_{\mathsf c(R)\in\mathcal G}
	 \bigl(h(R)-h(VR)-v(R)+v(HR)\bigr)=0.
	\end{equation}
	The first equality holds because every non-cut internal copy occurs in
	exactly two displayed incidence lists with opposite signs; the second
	holds because each parenthesis is $HVR-VHR=0$ after the canonical
	expansion proved above.  Applying $\mathcal P$ gives an equality in $C$;
	\eqref{eq:complete-ledger} separately retains the raw ancestry labels.

	We can now read the external boundary of the global trace cut as in
	Rules (2) and (3).  The sum of the evaluated initial rows is $dx=\Theta$;
	since every transfer and truncation is a literal identity, Rule (4)
	identifies its outside one-factor, weight-$(n+1)$ read with
	$\iota_C(\pr_{1,n+1}\Theta)$.  Equations \eqref{eq:indexed-horizontal-read} and
	\eqref{eq:indexed-vertical-read} identify the two oriented diagonal
	sums, and \eqref{eq:factor-two-occurrence-map} identifies the bottom
	sum.  The exhaustive terminal partition just proved makes every other
	external incidence silent.  The signs are read directly from
	\eqref{eq:four-incidences}: the homogeneous truncation child occurs with
	a minus sign in $h(R)=R-HR$, its vertical partner occurs with a plus sign
	in $v(HR)$, and the terminal factor-two incidence is on the negatively
	oriented factor boundary.  Taking the outside edge positively, the four, and
	only four, nonsilent families are therefore
	\begin{center}
	\begin{tabular}{c|c}
	 external family & contribution in $C$ \\ \hline
	 outside factor-one edge &
	 $\iota_C(\pr_{1,n+1}\Theta)$ \\
	 diagonal horizontal incidences, $2\leq k<n$ &
	 $-\Phi_k(f_k\chi_k)$ \\
	 oriented factor-two incidence & $-Q_n(\chi_n)$ \\
	 diagonal vertical incidences, $2\leq k<n$ &
	 $+T_k(\partial\chi_k)$
	\end{tabular}
	\end{center}
	Consequently
	\begin{equation}\label{eq:evaluated-complete-ledger}
	 0=\iota_C(\pr_{1,n+1}\Theta)
	   -\sum_{k=2}^{n-1}\Phi_k(f_k\chi_k)-Q_n(\chi_n)
	   +\sum_{k=2}^{n-1}T_k(\partial\chi_k).
	\end{equation}
	No homogeneous relation component has been called a relation in this
	argument: full relations occur only in the marked rows, whereas their
	homogeneous outputs are evaluated as elements of $F$.

	Finally, $T_k$ is supported only on
	$F_k\Sym^{m_k-1}(F_1)$, of total weight $n+1$.  The first assertion of
	the lemma identifies that component of $\partial\chi_k$ with the
	weight-$(n+1)$, factor-$m_k$ PBW coefficient of $\Theta$, which is zero.
	Thus $T_k(\partial\chi_k)=0$ for every $2\leq k<n$.  By
	\eqref{eq:terminal-definition}, equation
	\eqref{eq:evaluated-complete-ledger} says
	$\varepsilon+\sum_{k=2}^{n-1}T_k(\partial\chi_k)=0$; hence it proves both
	\eqref{eq:terminal-discrepancy} and $\varepsilon=0$.

	For $n=2$ the same construction has no indexed sets $\Delta_k$ and no
	diagonal incidences.  Its boundary therefore consists only of the
	outside factor-one incidence and the factor-two incidence, so
	\eqref{eq:evaluated-complete-ledger} reads
	\[
	 \iota_C(\pr_{1,3}\Theta)-Q_2(\chi_2)=0.
	\]
	Thus the empty-sum case is included without appealing to a separate
	pairing assertion.
	\end{proof}

\begin{lemma}[PBW assembly]\label{lem:PBW-assembly}
Let $\theta\in\Rel U$ satisfy \eqref{eq:governing}.  There are chains
$\chi_k$ as in \eqref{eq:packet-chains} such that, in
$C\subseteq\Q/\Z$,
\begin{equation}\label{eq:PBW-normal-form}
 \frac{\zeta}{q}+\Z
 =\sum_{k=2}^{n-1}\Phi_k(f_k\chi_k)
 -\eta_n\bigl(L_1\Sym^2(\pi_n)[\chi_n]\bigr).
\end{equation}
\end{lemma}

\begin{proof}
Apply Lemma~\ref{lem:closed-square} to $\Theta=\theta$.  By
\eqref{eq:terminal-definition} and the conclusion $\varepsilon=0$ of
that lemma,
\[
 \iota_C(\pr_{1,n+1}\theta)
 =\sum_{k=2}^{n-1}\Phi_k(f_k\chi_k)+Q_n(\chi_n).
\]
The occurrence-level saturation proof in that lemma is exactly what
identifies all three terms in this equality; no new frontier
classification is needed here.  Equation \eqref{eq:primitive-witness}
identifies the left side with $\zeta/q+\Z$.
Proposition~\ref{prop:quadratic-block} identifies the oriented quadratic term with
\[
 Q_n(\chi_n)
 =-\eta_n\bigl(L_1\Sym^2(\pi_n)[\chi_n]\bigr).
\]
Substitution gives \eqref{eq:PBW-normal-form}.
\end{proof}

\begin{proposition}[intermediate packet vanishing]
\label{prop:intermediate-vanishing}
For every $2\leq k<n$,
\[
                         \Phi_k(f_k\chi_k)=0.
\]
\end{proposition}

\begin{proof}
By \eqref{eq:chain-transgression},
$\Phi_k(f_k\chi_k)=T_k(\partial\chi_k)$.  The support calculation in
the proof of Lemma~\ref{lem:closed-square} and
\eqref{eq:governing} make the right-hand side zero.
\end{proof}

\section{Proof of the theorem}

\begin{proof}[Proof of Theorem~\ref{thm:main}]
The case $n=1$ was proved in Proposition~\ref{prop:top-reduction}, so
assume $n\geq2$.
By Proposition~\ref{prop:top-reduction}, let
$\zeta z\in\delta_{2n+1}(L)$ and choose the witness
\eqref{eq:witness}.  Lemma~\ref{lem:PBW-assembly} gives
\eqref{eq:PBW-normal-form}.  For every $k<n$,
Proposition~\ref{prop:intermediate-vanishing} makes the $k$th summand zero.
Proposition~\ref{prop:quadratic-vanishing} makes the last summand zero.
Hence
\[
                         \frac\zeta q+\Z=0,
\]
so $q\mid\zeta$ and $\zeta z=0$.
Proposition~\ref{prop:top-reduction} completes the induction.
\end{proof}

\begin{thebibliography}{9}

\bibitem{BM}
L.~Breen and R.~Mikhailov,
\emph{Derived functors of non-additive functors and homotopy theory},
Algebr. Geom. Topol. \textbf{11} (2011), 327--415;
\href{https://arxiv.org/abs/0910.2817}{arXiv:0910.2817}.

\bibitem{PS}
I.~B.~S.~Passi and T.~Sicking,
\emph{Dimension quotients of metabelian Lie rings},
Internat. J. Algebra Comput. \textbf{27} (2017), 251--258;
\href{https://arxiv.org/abs/1602.04919}{arXiv:1602.04919}.

\end{thebibliography}

\end{document}
